\documentclass{art}
\usepackage[utf8]{inputenc}
\usepackage[T1]{fontenc}
\usepackage{geometry}
\usepackage{graphicx}
\usepackage{float}
\usepackage{hyperref}
\usepackage[english,indonesian]{babel}

\title{Optimasi Pemrosesan Template LaTeX Menggunakan Pendekatan Berbasis AI}
\author{Academic Support Agent}
\date{\today}

\begin{document}
\maketitle

\begin{abstract}

\end{abstract}

\chapter{PENDAHULUAN}


\chapter{TINJAUAN PUSTAKA}


% Original content (commented out):
% % Optional: use minted for code listings
%\RequirePackage[]{minted}

\documentclass[
  a4paper,        % select between `a4paper' and `letterpaper'
  techreport,     % select between  `techreport', `report', `article', `commun', `persp' and `review'
  % lineno,         % switch, adds line numbering
  % easyreview,     % use easyReview package
  % anonymous,
  % noauthors,
  % nodate,
]{art}

\input{metadata.tex}
\input{macros.tex}
%\input{templates/pandoc.tex} % needed for pandoc-generated files

\begin{document}
\maketitle

\tableofcontents

\section{Introduction}

This document is an example of using various LaTeX features for writing ARTs.

\input{examples/test.tex}
%\input{examples/code.tex}

\section{Concluding remarks}

\section{Acknowledgements}

\nocite{*} % include all refs, even if not cited
\bibliography{art.bib,ref.bib}

\appendix

\section{Some extra stuff}


\end{document}


\end{document}